% Capitulo: Scrum com Kanban: o modelo que usei
\chapter{Scrum com Kanban: o modelo que usei}
\label{cap:scrumban}

\rascunho{
  Este capítulo foi montado a partir dos princípios do Scrum com Kanban e das ideias do
  livro. Revisar cada seção e substituir as escolhas genéricas pelas que eu de fato usei
  (duração das iterações, colunas do quadro, limites, ritos), marcadas como pendências.
}

O capítulo anterior mostrou por que o Scrum aplicado ao pé da letra costuma frustrar times de
dados, e o que vale aproveitar dele e do Kanban. Este capítulo descreve um modelo que combina
os dois, aplicado em projetos de visualização e de engenharia de dados. Ele foi pensado para
ser copiado e adaptado: nenhuma das escolhas abaixo é a única certa, mas cada uma resolve um
problema concreto, e vale entender o problema antes de mudar a escolha \parencite{ladas2009}.

\section{Cadências}
\label{sec:cadencias}

Do Scrum, o modelo mantém as cadências, porque são elas que mantêm o negócio por perto e o
time melhorando. O que muda é o papel delas: em vez de marcar um compromisso fechado de
escopo, cada cadência é um momento de decisão e de ajuste.

O \textbf{planejamento} acontece no início de cada iteração \pendente{duração das iterações
no meu modelo}. Nele, o time e quem prioriza o trabalho revisam o que está pronto para entrar
e escolhem o que puxar primeiro, respeitando os limites de trabalho em andamento. Não há
compromisso de terminar uma lista fechada; há uma ordem de prioridade e uma previsão baseada
no histórico do fluxo (\veja{sec:metricas-fluxo}).

O \textbf{alinhamento diário} é curto e olha para o quadro, da direita para a esquerda: o que
está perto de ficar pronto e o que está travado, antes do que vai começar. A pergunta central
não é \enquote{o que cada um fez ontem}, e sim \enquote{o que impede o trabalho de fluir}.

A \textbf{revisão} mostra a quem pediu o que ficou pronto, com dados reais, e é o momento de
validar e de corrigir o rumo. É também o lugar natural para a revisão periódica dos produtos
de dados em uso, com as métricas de acesso e de custo do capítulo~\ref{cap:frequencia}.

A \textbf{retrospectiva} discute o processo: o que travou, que regra do quadro precisa mudar,
que tipo de demanda está chegando mal definida. É o rito que faz o modelo evoluir, e o
primeiro a ser cortado quando o time está sobrecarregado, o que é justamente quando ele mais
faz falta.

\rascunho{
  Confirmar quais desses ritos existiam no meu modelo, com que frequência e duração, e se
  havia algum rito próprio (por exemplo, uma reunião de refinamento com o negócio).
}

\section{O quadro e o limite de WIP}
\label{sec:quadro-wip}

Do Kanban, o modelo traz o quadro e a gestão do fluxo \parencite{anderson2010}. A
Figura~\ref{fig:quadro-scrumban} mostra um exemplo de quadro para um time de dados.

\begin{figure}[htbp]
  \centering
  % Exemplo de quadro de Scrum com Kanban para um time de dados (cap. 28).
\begin{tikzpicture}[diagrama, x=1mm, y=1mm,
    cab/.style={font=\scriptsize\bfseries, align=center, text width=14mm},
    raia/.style={font=\tiny, text=cinza, align=left, text width=13mm, anchor=west},
    cartao/.style={draw=tinta, fill=white, minimum width=11mm, minimum height=3.2mm, inner sep=0pt},
    cartaourg/.style={cartao, fill=fundo, line width=0.7pt}]
  % colunas
  \foreach \i/\nome in {0/Entrada, 1/Descoberta, 2/{Pronto p/\\fazer}, 3/{Fazendo\\(até 3)},
                        4/{Revisão\\(até 2)}, 5/{Validação\\negócio}, 6/Pronto} {
    \node[cab] at (21.5+\i*15,-3) {\nome};
  }
  % grade
  \draw[fio] (14,-7.5) -- (119,-7.5);
  \foreach \i in {0,...,7} \draw[fio] (14+\i*15,0) -- (14+\i*15,-49);
  \foreach \y in {-21.5,-35.5,-49} \draw[fio] (0,\y) -- (119,\y);
  % raias
  \node[raia] at (0,-14.5) {Urgente};
  \node[raia] at (0,-28.5) {Projeto};
  \node[raia] at (0,-42.5) {Sustentação};
  % cartoes (exemplo)
  \node[cartaourg] at (66.5,-14.5) {};
  \foreach \c/\y in {0/-25.5, 0/-30.5, 1/-28, 2/-25.5, 2/-30.5, 3/-25.5, 3/-30.5, 4/-28, 5/-28, 6/-25.5, 6/-30.5}
    \node[cartao] at (21.5+\c*15,\y) {};
  \foreach \c/\y in {0/-40, 0/-45, 3/-42.5, 6/-40, 6/-45}
    \node[cartao] at (21.5+\c*15,\y) {};
\end{tikzpicture}

  \caption{Um exemplo de quadro de Scrum com Kanban para um time de dados, com colunas que
    separam descoberta de construção e raias por classe de serviço. O cartão destacado é
    urgente.}
  \label{fig:quadro-scrumban}
  \fonte{elaborado pelo autor.}
\end{figure}

Três escolhas do quadro merecem explicação. A primeira é a coluna de \textbf{descoberta},
separada da construção. É onde ficam as investigações curtas que olham os dados antes de o
trabalho ser refinado (\veja{sec:descoberta}). Um item só passa para \enquote{pronto para
fazer} quando as perguntas que definem o tamanho dele foram respondidas.

A segunda é a coluna de \textbf{validação com o negócio}, separada da revisão técnica. A
revisão confere se o código está certo; a validação confere se o resultado responde à pergunta
de quem pediu. São coisas diferentes, feitas por pessoas diferentes, e misturá-las numa coluna
só esconde onde o trabalho está parado.

A terceira são as \textbf{raias por classe de serviço}. Nem todo trabalho tem a mesma
urgência, e o Kanban propõe tratar cada tipo de forma explícita: itens urgentes, que passam na
frente e têm limite de quantidade; itens de projeto, que seguem a prioridade normal; e itens
de sustentação, que precisam de uma parte reservada da capacidade para não competir com os
projetos a cada dia \parencite{anderson2010}. Visualizar a sustentação no quadro, em vez de
tratá-la como interrupção invisível, é uma das mudanças que mais ajudam um time de dados.

Os números ao lado das colunas \enquote{fazendo} e \enquote{revisão} são os limites de
trabalho em andamento. Eles parecem uma restrição, e são. A razão de existirem vem da teoria
das filas: num sistema estável, o tempo médio que um item leva para atravessar o sistema é
igual à quantidade média de itens em andamento dividida pela vazão
\parencite{little1961,reinertsen2009}. Se o time termina, em média, cinco itens por semana, ter
vinte itens em andamento significa que cada um leva, em média, quatro semanas. Limitar o
trabalho em andamento é a forma mais direta de encurtar o tempo de cada entrega, e de fazer o
time terminar antes de começar.

\rascunho{
  Confirmar as colunas e os limites do meu quadro real, a ferramenta usada e se havia raias
  ou classes de serviço diferentes das do exemplo.
}

\section{Refinamento}
\label{sec:refinamento}

Boa parte dos problemas de um time de dados começa antes do trabalho: pedidos vagos que viram
histórias vagas. O refinamento é o momento de transformar um pedido em algo que o time
consegue construir e o negócio consegue validar, e vale ter critérios explícitos para que uma
história possa entrar na coluna \enquote{pronto para fazer}.

Para histórias de dados, alguns elementos fazem muita diferença. A pergunta de negócio, ou a
decisão que a entrega vai apoiar, escrita de forma que alguém de fora entenda. A forma
tradicional das histórias de usuário ajuda: \enquote{como gerente comercial, quero ver a
margem por vendedor para decidir a revisão das metas} diz mais do que \enquote{painel de
margem} \parencite{cohn2004}. As fontes envolvidas, já investigadas na descoberta. O dono de
negócio de cada métrica. A frequência de atualização necessária, definida a partir do ritmo
da decisão, e não do pedido (\veja{sec:ritmo-decisao}).

E, cada vez mais, os testes de aceite. Com a IA, escrever os testes ficou barato o bastante
para que eles sejam escritos já no refinamento, junto com o stakeholder, a partir das regras
de negócio que ele descreve (\veja{sec:tdd}). A história entra no quadro com a regra
transformada em teste, e o time passa a construir para fazer o teste passar, em vez de
construir e depois descobrir, na validação, que a regra era outra.

\section{Definição de pronto}
\label{sec:definicao-pronto}

A definição de pronto é a lista do que precisa estar feito para um item sair da última coluna.
Ela protege o time da entrega que parece terminada e não está, e protege o negócio do painel
que funciona hoje e quebra na semana que vem. Para dados, uma definição de pronto útil
costuma incluir:

\begin{itemize}
  \item o código revisado por outra pessoa (\veja{sec:revisao-codigo});
  \item os testes escritos antes, e passando (\veja{sec:tdd});
  \item as tabelas, colunas e métricas novas documentadas;
  \item a frequência de atualização combinada e monitorada;
  \item o painel mostrando a data da última atualização;
  \item a publicação feita pelo processo automático, e não à mão (capítulo~\ref{cap:ci-cd});
  \item a validação de quem pediu, com dados reais;
  \item um dono definido para o produto entregue.
\end{itemize}

A lista parece longa, mas a maior parte dela é automática quando as práticas da
Parte~\ref{parte:devops} estão no lugar. E é ela que garante que cada entrega deixe o
ambiente de dados um pouco melhor, e não um pouco mais frágil.

\section{Métricas do fluxo}
\label{sec:metricas-fluxo}

O modelo troca as estimativas detalhadas por medição do fluxo real. Quatro indicadores
bastam para a maior parte das conversas. O \textbf{tempo de ciclo}: quanto tempo um item leva
do momento em que o trabalho começa até ficar pronto. A \textbf{vazão}: quantos itens ficam
prontos por semana. O \textbf{trabalho em andamento}: quantos itens estão sendo trabalhados ao
mesmo tempo. E a \textbf{idade dos itens em andamento}: há quanto tempo cada item está no
quadro, o que revela cedo o que está travado.

Com o histórico desses indicadores, as previsões ficam mais honestas do que qualquer
estimativa em pontos. Em vez de \enquote{essa história vale oito pontos}, o time pode dizer
\enquote{itens desse tipo levam, em 85\% dos casos, até nove dias}. É uma resposta que o
negócio entende e que incorpora, sem esforço adicional, as dependências, a sustentação e a
descoberta que as estimativas costumam ignorar.

\rascunho{
  Relato: como esse modelo foi implantado (o que mudou em relação ao Scrum anterior, como o
  time e o negócio reagiram, o que melhorou nos indicadores). Números só em ordem de grandeza.
}

\begin{leituras}
  \leitura{ladas2009}{O livro que propôs o Scrumban como caminho de evolução do Scrum para um
    sistema de fluxo.}
  \leitura{reinertsen2009}{Para entender, com a teoria das filas, por que limitar o trabalho
    em andamento acelera as entregas.}
  \leitura{cohn2004}{A referência sobre histórias de usuário, útil para escrever histórias de
    dados que o negócio entende.}
\end{leituras}
